\documentclass[../../main.tex]{subfiles}

\begin{document}

\chapter[Orientasi Buku Ajar]{Orientasi: Buku Ajar, OBE, dan Peta Perjalanan}
\label{chap:orientasi}

\begin{outcomebox}
\textbf{CPMK yang menjadi konteks:} \cpmk{01}, \cpmk{02}, \cpmk{03}, \cpmk{04}\\[3pt]
Setelah menyelesaikan bab ini, mahasiswa mampu:
\begin{enumerate}[label=\arabic*.,leftmargin=*]
  \item menjelaskan tujuan, cakupan, dan struktur buku ajar;
  \item menghubungkan setiap bab dengan RPS berbasis OBE;
  \item menggunakan fitur pembelajaran, proyek mini-kompilator, dan
        sumber belajar yang tersedia;
  \item menyusun rencana belajar pribadi untuk satu semester.
\end{enumerate}
\end{outcomebox}

\section*{Materi Pokok}
\begin{itemize}[leftmargin=*]
  \item Mengapa Teknik Kompilasi penting bagi sarjana informatika.
  \item Posisi mata kuliah dalam kurikulum dan RPS berbasis OBE.
  \item Arsitektur buku dan alur belajar tujuh belas bab.
  \item Anatomi bab, kotak pembelajaran, dan komponen asesmen.
  \item Proyek mini-kompilator Subset-C dan strategi belajar.
\end{itemize}

\subfile{section-01}
\subfile{section-02}
\subfile{section-03}
\subfile{section-04}
\subfile{section-05}

\section*{Contoh}
\addcontentsline{toc}{section}{Contoh}

\begin{contoh}[Membaca anatomi satu bab]
Perhatikan sebuah bab inti, misalnya Bab VII tentang implementasi
lexer. Kotak outcome menyebut Sub-CPMK 2.4, yaitu mengimplementasikan
lexical analyzer untuk bahasa sederhana. Bagian materi menjelaskan
buffer, sentinel, dan \eng{maximal munch}. Contoh menunjukkan potongan
kode \code{next\_token}. Praktikum meminta mahasiswa melengkapi lexer
dan menjalankan 30 tes. Asesmen menilai kebenaran, kualitas kode,
kelengkapan, dan pengujian. Daftar cek menyatakan bukti yang harus
ditunjukkan.

Dengan membaca anatomi ini, mahasiswa mengetahui sejak awal apa yang
harus dihasilkan, bukan hanya apa yang harus dibaca.
\end{contoh}

\begin{contoh}[Dari program sumber ke token]
Untuk program pada \lstref{lst:ch01-hello}, urutan token yang dihasilkan
lexer kira-kira adalah:
\begin{center}
\code{int}, \code{main}, \code{(}, \code{)}, \code{\{}, \code{int},
\code{a}, \code{=}, \code{5}, \code{;}, \code{int}, \code{b},
\code{=}, \code{7}, \code{;}, \code{print}, \code{(}, \code{a},
\code{+}, \code{b}, \code{)}, \code{;}, \code{return}, \code{0},
\code{;}, \code{\}}, \code{EOF}.
\end{center}
Perhatikan bahwa spasi dan komentar tidak menghasilkan token, sedangkan
setiap tanda baca dan literal menghasilkan token tersendiri.
\end{contoh}

\section*{Aktivitas Pembelajaran dan Praktikum}
\addcontentsline{toc}{section}{Aktivitas Pembelajaran dan Praktikum}

\begin{table}[H]
\centering
\caption{Rencana aktivitas Bab I}
\label{tab:ch01-aktivitas}
\begin{tabularx}{\textwidth}{@{}llX@{}}
\toprule
\textbf{Sesi} & \textbf{Bentuk} & \textbf{Kegiatan} \\
\midrule
1 & Ceramah & Membaca peta buku dan RPS; diskusi capaian pembelajaran \\
2 & Latihan & Menyusun rencana belajar pribadi dan memeriksa prasyarat \\
3 & Praktikum 0 & Menyiapkan lingkungan kerja dan menjalankan program contoh \\
\bottomrule
\end{tabularx}
\end{table}

\subsection*{Praktikum 0: Persiapan Lingkungan}
\begin{enumerate}[label=\arabic*.,leftmargin=*]
  \item Pastikan \code{pdflatex} dan \code{bibtex} tersedia.
  \item Unduh atau salin direktori proyek \file{code/subsetc}.
  \item Baca \file{code/subsetc/README.md} dan spesifikasi bahasa.
  \item Jalankan program contoh \file{hello.sc} melalui alat yang
        disediakan dosen.
  \item Tuliskan satu paragraf refleksi tentang kesulitan yang dihadapi.
\end{enumerate}

Praktikum ini tidak dinilai secara sumatif, tetapi menjadi prasyarat
administratif untuk Praktikum 1.

\section*{Latihan dan Refleksi}
\addcontentsline{toc}{section}{Latihan dan Refleksi}

\subsection*{Latihan}
\begin{enumerate}[label=\arabic*.,leftmargin=*]
  \item Jelaskan perbedaan tujuan buku ajar ini dengan buku teks
        kompilator pada umumnya.
  \item Sebutkan empat CPMK dan hubungkan masing-masing dengan minimal
        satu bab inti.
  \item Mengapa buku ini memisahkan bab orientasi dan landasan teori
        dari empat belas bab inti?
  \item Gambarkan alur dari program sumber sampai AST dalam bentuk
        diagram sendiri.
  \item Baca \file{code/subsetc/examples/hello.sc}, lalu tuliskan
        token-token yang dihasilkan secara manual.
\end{enumerate}

\subsection*{Refleksi}
\begin{reflectionbox}
Jawab dengan jujur:
\begin{enumerate}[label=\arabic*.,leftmargin=*]
  \item Kompetensi apa yang paling ingin Anda capai pada mata kuliah ini?
  \item Bagian mana dari proses kompilasi yang paling sulit Anda bayangkan?
  \item Kebiasaan belajar apa yang akan Anda ubah setelah membaca bab ini?
\end{enumerate}
\end{reflectionbox}

\section*{Asesmen}
\addcontentsline{toc}{section}{Asesmen}

Asesmen Bab I bersifat formatif dan tidak memiliki bobot tersendiri.
Nilainya menjadi bagian dari komponen \textbf{Kuis} (5\%) dan
\textbf{Tugas Teori} (8\%) pada RPS.

\begin{table}[H]
\centering
\caption{Instrumen asesmen dan indikator}
\label{tab:ch01-asesmen}
\begin{tabularx}{\textwidth}{@{}lXl@{}}
\toprule
\textbf{Instrumen} & \textbf{Indikator} & \textbf{Bobot} \\
\midrule
Kuis awal & Menjelaskan struktur buku, CPMK, dan alur belajar &
5\% (bagian) \\
Tugas rencana belajar & Menyusun rencana studi yang realistis dan
terukur & 8\% (bagian) \\
Praktikum 0 & Menyiapkan lingkungan dan menjalankan program contoh &
kelulusan prasyarat \\
\bottomrule
\end{tabularx}
\end{table}

Kriteria kelulusan Bab I: mahasiswa dapat menjelaskan peta buku,
menghubungkan bab dengan CPMK, dan menunjukkan bahwa lingkungan kerja
siap dipakai.

\section*{Checklist Pencapaian Kompetensi}
\addcontentsline{toc}{section}{Checklist Pencapaian Kompetensi}

Centang setiap butir yang sudah dapat Anda tunjukkan dengan bukti.

\begin{checklist}
  \item Saya dapat menjelaskan tujuan dan cakupan buku ajar.
  \item Saya dapat menyebutkan empat CPMK dan kaitannya dengan bab.
  \item Saya dapat menjelaskan alur tujuh belas bab.
  \item Saya dapat menemukan komponen OBE di setiap bab.
  \item Saya sudah menyiapkan lingkungan kerja untuk praktikum.
  \item Saya sudah menjalankan program contoh \file{hello.sc}.
  \item Saya sudah menuliskan rencana belajar pribadi.
  \item Saya dapat menjelaskan perbedaan front-end dan back-end.
\end{checklist}

\section*{Rangkuman}
\addcontentsline{toc}{section}{Rangkuman}

Pada bab ini telah dibahas:
\begin{itemize}[leftmargin=*]
  \item pentingnya Teknik Kompilasi bagi sarjana informatika;
  \item posisi mata kuliah dalam kurikulum dan RPS berbasis OBE;
  \item arsitektur buku yang terdiri atas orientasi, landasan teori,
        empat belas bab inti, dan bab evaluasi;
  \item anatomi bab serta fungsi kotak pembelajaran dan kode program;
  \item proyek mini-kompilator Subset-C dan milestone-nya.
\end{itemize}

Bab berikutnya, Bab II, menyajikan landasan teori dan notasi formal
yang dipakai di seluruh buku: himpunan, relasi, bahasa, tata bahasa,
dan mesin abstrak. Materi tersebut menjadi bahasa bersama untuk
memahami analisis leksikal, sintaksis, dan semantik.

\section*{Evaluasi}
\addcontentsline{toc}{section}{Evaluasi}

\begin{enumerate}[label=\arabic*.,leftmargin=*]
  \item Uraikan tiga alasan mengapa pemahaman teknik kompilasi berguna
        bagi seorang insinyur perangkat lunak.
  \item Buat tabel yang memetakan empat CPMK ke minimal dua bab inti
        masing-masing.
  \item Jelaskan perbedaan peran Bab I, Bab II, dan Bab XVII dalam
        struktur buku.
  \item Diberikan potongan program Subset-C, tuliskan urutan token dan
        jelaskan aturan yang Anda pakai.
  \item Rancang rencana belajar empat belas minggu yang memuat bacaan,
        praktikum, dan evaluasi mandiri.
\end{enumerate}

\subsection*{Kunci dan Rubrik Singkat}
\begin{itemize}[leftmargin=*]
  \item Soal 1--3: ketepatan konsep (60\%), kelengkapan (25\%),
        kejelasan bahasa (15\%).
  \item Soal 4: kebenaran token (70\%), penjelasan aturan (30\%).
  \item Soal 5: kelayakan jadwal (50\%), keterukuran target (30\%),
        keseimbangan teori-praktik (20\%).
\end{itemize}


\ifSubfilesClassLoaded{%
  \bibliographystyle{apalike}%
  \bibliography{../../references}%
}{}

\end{document}
